\chapter{Introduction}

\section{Background}
Lung disease remains a leading cause of global morbidity and mortality. In the diagnosis of lung disease, assessing the pattern,
location, and regional distribution of involvement is the domain of radiology. Clinical diagnosis increasingly relies on
radiological patterns to differentiate between pathologies with overlapping symptoms. While Chest X-rays (CXR) remain the first-line
screening tool due to low cost and radiation, they suffer from limited sensitivity (roughly 70\% visibility of lung volume due to
anatomical superposition, in which all anatomical structures located in the path of the X-rays are stacked). [1]

CT is a radiographic technology that combines X-ray equipment with a computer and a cathode ray tube display to produce images
of cross sections of the human body. [2] **(Advantage of CT vs X-ray).
The fundamental metric of a CT scan is the Hounsfield Unit (HU), a relative quantitative scale that describes radiodensity [1.1]
In pulmonary imaging, HU values allow for the objective differentiation of tissues that might appear visually similar.

A key characteristic in CT is "ground-glass opacity" (GGO), which appears as hazy increased opacity of lung, with preservation
of bronchial and vascular margins. It is caused by partial filling of airspaces, interstitial thickening, collapse of alveoli,
and increased capillary blood volume.
"Consolidation" on the other hand, appears as a homogeneous increase in pulmonary parenchymal attenuation that obscures the margins
of vessels and airway walls (45) (Fig 19). An air bronchogram may be present. The attenuation characteristics of consolidated
lung are only rarely helpful in differential diagnosis (eg, decreased attenuation in lipoid pneumonia [46] and increased in amiodaron
toxicity. Consolidation is more opaque than \ac{GGO}. While \ac{GGO} often suggests a reversible or early-stage process, consolidation
usually indicates a more aggressive or advanced stage. The Consolidation-to-GGO (C/G) ratio is a quantitative metric calculated by dividing
the volume or area of consolidation by the volume or area of ground-glass opacity within a lesion or throughout the lungs.

The traditional approach for CT analysis relies on a "slice-by-slice" approach where a radiologist views a series of axial
cross-sections and identifies the "index slice"—the single image where a pathology appears at its largest diameter.
A 2D slice only captures a small fraction of a 3D structure. If a nodule is irregularly shaped, the "largest" slice may not
represent the true extent of the disease [2.2]. Studies show that 2D measurements are highly sensitive to the patient’s
orientation and the specific slice chosen, leading to high inter-observer variability [2.3]. 2D models treat each slice as an
independent data point, losing the "spatial continuity" between adjacent tissues [1.4].
3D volumetric analysis represents a quantitative shift wherein the entire dataset is used to reconstruct the lung's anatomy as a
solid volume. 3D advantages?? [5.1].

Modern radiology increasingly relies on Artificial Intelligence (AI) to streamline the acquisition of radiologist analyses on
chest X-rays, as evidenced by a study wherein an \ac{AI} system reduced interpretation delivery times from 11.2 days to a mere
2.7 days, reinforcing the potency of automated triaging systems in streamlining healthcare workflows and amplifying patient 
care standards [42].[1] However, the exponential increase in data volume generated by 3D CT scanners presents a profound challenge
for human interpretation and computational analysis. While the prevailing research trend has focused heavily on the deployment of
3D deep learning architectures, such as Convolutional Neural Networks (CNNs), these models frequently encounter significant hurdles
regarding computational overhead, data scarcity, and a critical lack of interpretability. The research proposed herein advocates
for a departure from these "black box" paradigms, suggesting that sparse representation methods offer a superior framework for the
volumetric identification of lung diseases by ensuring clinical interpretability.

\section{Problem Statement}
Modern radiology increasingly relies on artificial intelligence to manage the vast quantities of 3D data produced during routine
screenings. 3D deep learning has been widely applied to the automatic diagnosis of lung diseases, achieving remarkable success in
tasks such as segmenting suspected infection regions and differentiating between various types of carcinoma. Despite these advancements,
the inherent complexity of 3D \ac{CNN}s introduces a set of constraints that are often overlooked in the pursuit of high accuracy scores.

\subsection{Computational Overhead}
Processing a 3D CT scan involves treating stacked 2D slices as a single volumetric matrix, which demands significantly more GPU
memory and processing power than standard image analysis. Unlike ordinary RGB images, CT scans consist of single-channel data with
varying voxel dimensions and resolution, which complicates spatial alignment and generalizability across different scanners.
Furthermore, the "curse of dimensionality" remains a persistent threat; as the network depth increases to capture global features,
the risk of gradient explosion and slow convergence becomes more pronounced.

\subsection{Interpretability Bottleneck}
Beyond computational limits, the clinical utility of deep learning is hampered by the interpretability paradox. While tools like
Grad-CAM provide visual heatmaps to localize lesions, these post-hoc explanations are often insufficient for medical professionals
who require a granular understanding of the model's decision-making process. In high-stakes environments, a model that shruggingly
produces a diagnosis without an audit trail of its reasoning is difficult to integrate into standard workflows, as it lacks the
accountability necessary for legal and ethical compliance. This transparency gap has led to a renewed interest in "open book" models
that are interpretable by design, where the internal logic is clear from the outset.

\section{Objectives and Scope}
\subsection{Objectives}
\begin{enumerate}
    \item Develop a Dictionary Learning framework for Interpretable Lung Disease Classification from 3D Chest CT
        \begin{itemize}
            \item Implement training in two phases for normal and pathological atoms
            \item Integrate \ac{KSVD} with label consistency constraints
        \end{itemize}
    \item Visualize dictionary atoms and map them to anatomical features
        \begin{itemize}
            \item Compare sparse coefficient patterns across disease classes
            \item Demonstrate atom-to-pathology correspondence
        \end{itemize}
    \item Benchmark against baseline models
        \begin{itemize}
            \item Compare with standard \ac{KSVD} without label consistency constraint
            \item Compare with 3D \ac{CNN} baseline (3D-ResNet)
        \end{itemize}
    \item UI ??
    \item Integrate clinical metrics for volumetric analysis
        \begin{itemize}
            \item Implement HU-based segmentation for GGO and consolidation
            \item Calculate C/G ratio and compare model predictions against clinical thresholds
        \end{itemize}
\end{enumerate}

\subsection{Scope}


% [1] Radiological Diagnosis in Lung Disease
% [2] A Comparative Study of Medical Imaging Techniques
% [3] A Review of Artificial Intelligence Integration in Medical Imaging